% ============================================================================
% LANGENTWURF - Physik Klasse 6: s(t)-Diagramme
% ============================================================================
% Tier 1 (FULL) | Profil A (Gesamtschule heterogen)
% Stunde 3 von 4 der Reihe "Bewegungen"
% ============================================================================

\documentclass[11pt,a4paper]{article}

% ============================================================================
% PACKAGES
% ============================================================================
\usepackage[utf8]{inputenc}
\usepackage[T1]{fontenc}
\usepackage[ngerman]{babel}
\usepackage{geometry}
\usepackage{fancyhdr}
\usepackage{lastpage}
\usepackage{xcolor}
\usepackage{tcolorbox}
\usepackage{enumitem}
\usepackage{tabularx}
\usepackage{longtable}
\usepackage{array}
\usepackage{booktabs}
\usepackage{hyperref}
\usepackage{ifthen}
\usepackage{tikz}
\usepackage{amssymb}
\usepackage{amsmath}
\usepackage{siunitx}
\sisetup{locale=DE}

% ============================================================================
% KONFIGURATION
% ============================================================================

\newcommand{\plantier}{1}
\newcommand{\lerngruppenprofil}{A}
\newcommand{\fachid}{physik}

\newcommand{\tierone}[1]{\ifthenelse{\equal{\plantier}{1}}{#1}{}}
\newcommand{\tieronetwo}[1]{\ifthenelse{\equal{\plantier}{1}\OR\equal{\plantier}{2}}{#1}{}}
\newcommand{\tierall}[1]{#1}

\newcommand{\testdifferenziert}[1]{%
    \ifthenelse{\equal{\fachid}{physik}\OR\equal{\fachid}{chemie}}{#1}{}%
}
\newcommand{\profilheterogen}[1]{%
    \ifthenelse{\equal{\lerngruppenprofil}{A}}{#1}{}%
}

% ============================================================================
% VARIABLEN
% ============================================================================

\newcommand{\fach}{Physik}
\newcommand{\jahrgangsstufe}{6}
\newcommand{\schulart}{Gesamtschule}
\newcommand{\kursart}{Orientierungsstufe}
\newcommand{\stundenthema}{Weg-Zeit-Diagramme verstehen und zeichnen}
\newcommand{\reihenthema}{Bewegungen -- Geradlinige, gleichförmige Bewegung}
\newcommand{\stundennummer}{3}
\newcommand{\reihenstunden}{4}
\newcommand{\unterrichtsdatum}{}
\newcommand{\unterrichtszeit}{}
\newcommand{\raum}{}

\newcommand{\lehrkraft}{J. Cuno}
\newcommand{\ausbildungsstand}{Lehrkraft}

\newcommand{\lerngruppengroesse}{24}
\newcommand{\maennlich}{12}
\newcommand{\weiblich}{12}
\newcommand{\divers}{0}

\newcommand{\fachfarbe}{2c5aa0}

% ============================================================================
% LAYOUT
% ============================================================================
\geometry{left=2.5cm, right=2.5cm, top=2.5cm, bottom=2.5cm}
\definecolor{fach}{HTML}{\fachfarbe}
\definecolor{gray}{HTML}{666666}
\definecolor{lightgray}{HTML}{f5f5f5}
\definecolor{successgreen}{HTML}{2a7a4b}
\definecolor{warningorange}{HTML}{b35c00}

\tcbuselibrary{skins,breakable}

\newtcolorbox{infobox}[1][]{
    colback=lightgray,
    colframe=fach,
    fonttitle=\bfseries,
    title=#1,
    breakable
}

\newtcolorbox{scorebox}{
    colback=white,
    colframe=fach,
    fonttitle=\bfseries,
    title=Didaktik-Score,
    breakable
}

\newtcolorbox{profilbox}{
    colback=lightgray,
    colframe=gray,
    fonttitle=\bfseries,
    title=Lerngruppen-Profil,
    breakable
}

\pagestyle{fancy}
\fancyhf{}
\fancyhead[L]{\small\textcolor{gray}{\fach{} -- Jg. \jahrgangsstufe{} -- \stundenthema}}
\fancyhead[R]{\small\textcolor{gray}{Langentwurf Tier \plantier{} / Profil \lerngruppenprofil}}
\fancyfoot[C]{\small\thepage{} / \pageref{LastPage}}
\renewcommand{\headrulewidth}{0.4pt}
\renewcommand{\footrulewidth}{0pt}

\renewcommand{\thesection}{\arabic{section}}
\renewcommand{\thesubsection}{\thesection.\arabic{subsection}}

% ============================================================================
% DOKUMENT
% ============================================================================
\begin{document}

% ============================================================================
% DECKBLATT
% ============================================================================
\begin{titlepage}
    \centering
    \vspace*{2cm}

    {\Large\textcolor{gray}{}}\\[0.5cm]
    {\large\textcolor{gray}{\schulart{} -- Orientierungsstufe}}\\[2cm]

    {\Huge\textcolor{fach}{\textbf{Langentwurf}}}\\[0.5cm]
    {\large Tier 1 -- Vollständig \quad|\quad Profil A}\\[2cm]

    {\LARGE\textbf{\stundenthema}}\\[0.5cm]
    {\large im Rahmen der Unterrichtsreihe}\\[0.3cm]
    {\Large\textit{\reihenthema}}\\[2cm]

    \begin{tabular}{rl}
        \textbf{Fach:} & \fach \\[0.3cm]
        \textbf{Klasse:} & \jahrgangsstufe{} (\kursart) \\[0.3cm]
        \textbf{Stunde:} & \stundennummer{} von \reihenstunden \\[0.3cm]
        \textbf{Datum:} & \underline{\hspace{3cm}} \\[0.3cm]
        \textbf{Zeit:} & \underline{\hspace{3cm}} \\[0.3cm]
        \textbf{Raum:} & \underline{\hspace{3cm}} \\[0.3cm]
    \end{tabular}

    \vfill

    {\large Lehrkraft: \lehrkraft}\\[0.3cm]
    {\normalsize\textcolor{gray}{\ausbildungsstand}}\\[1cm]

\end{titlepage}

% ============================================================================
% INHALTSVERZEICHNIS
% ============================================================================
\tableofcontents
\newpage

% ============================================================================
% KONFIGURATIONS-ÜBERSICHT
% ============================================================================
\begin{infobox}[Konfiguration dieses Langentwurfs]
\begin{tabular}{ll}
    \textbf{Tier:} & 1 -- Vollständig (alle Abschnitte) \\[0.3em]
    \textbf{Profil:} & A -- Gesamtschule heterogen \\[0.3em]
    \textbf{Fach:} & Physik \\[0.3em]
    \textbf{Testdifferenzierung:} & Ja (3 Niveaus: */~**/~***)
\end{tabular}
\end{infobox}

% ============================================================================
% LERNGRUPPEN-PROFIL
% ============================================================================
\begin{profilbox}
\textbf{Profil A: Gesamtschule heterogen}\\[0.5em]
\begin{tabular}{ll}
    Klassenstufe: & 6 (Orientierungsstufe) \\
    Zusammensetzung: & BR (20\%), MR (50\%), GY (30\%) \\
    Inklusion: & 2--3 SuS mit LRS, ggf. 1 SuS mit Förderbedarf \\
    Testdifferenzierung: & 3 Niveaus (*/**/***) \\
    Besonderheiten: & Hohe Heterogenität, intensive Differenzierung nötig
\end{tabular}
\end{profilbox}

% ============================================================================
% 1. BEDINGUNGSANALYSE
% ============================================================================
\section{Bedingungsanalyse}

\subsection{Statistik der Lerngruppe}
\begin{tabular}{ll}
    Klassengröße: & \lerngruppengroesse{} SuS (\maennlich{} m / \weiblich{} w / \divers{} d) \\
    Jahrgangsstufe: & \jahrgangsstufe \\
    Kursart: & \kursart \\
    Profil: & \lerngruppenprofil{} (heterogen) \\
\end{tabular}

\subsection{Zusammensetzung nach Leistungsniveau}
\begin{tabular}{|l|c|l|}
\hline
\textbf{Niveau} & \textbf{Anteil} & \textbf{Charakteristik} \\
\hline
Basis (*) & ca. 20\% & Qualitative Beschreibungen, Ablesen, Zuordnen \\
\hline
Standard (**) & ca. 50\% & Zeichnen von Diagrammen, einfache Berechnungen \\
\hline
Erweitert (***) & ca. 30\% & Interpretation, Vergleich, Transfer \\
\hline
\end{tabular}

\subsection{Individuelle Besonderheiten (Inklusion)}
\begin{itemize}
    \item \textbf{LRS:} 2--3 SuS (Nachteilsausgleich: Zeitzugabe, vergrößerte Schrift, vorstrukturierte Diagramme)
    \item \textbf{Visuell-räumliche Schwächen:} Ggf. 1--2 SuS (zusätzliche Unterstützung beim Zeichnen, Raster vorgeben)
    \item \textbf{Besonders leistungsstark:} 2--3 SuS (Zusatzaufgaben: eigene Bewegungsgeschichten)
    \item \textbf{DaZ:} Ggf. 1--2 SuS (Wortschatzarbeit: Achsenbeschriftung, Fachbegriffe)
\end{itemize}

\subsection{Arbeitsvoraussetzungen}
\begin{itemize}
    \item Raumsituation: Physikfachraum oder Klassenraum mit Tafel/Whiteboard
    \item Technische Ausstattung: Beamer, WLAN, Tablets/PCs (optional für Simulation)
    \item Verfügbare Medien: Tafel, Arbeitsblätter, digitaler Lernpfad, Millimeterpapier
\end{itemize}

\subsection{Fachbezogener Entwicklungsstand}
\begin{itemize}
    \item \textbf{Vorwissen (Kl. 6 -- bisherige Stunden):}
    \begin{itemize}
        \item Stunde 1: Bewegungsarten (Bahnform: geradlinig, kreisförmig, krummlinig; Bewegungsart: gleichförmig, ungleichförmig)
        \item Stunde 2: Geschwindigkeit messen (Weg $s$, Zeit $t$, Formel $v = \frac{s}{t}$, Einheiten m/s, km/h)
    \end{itemize}
    \item \textbf{Mathematisches Vorwissen (Kl. 5/6):}
    \begin{itemize}
        \item Koordinatensystem (x-/y-Achse) aus Mathematik bekannt
        \item Punkte einzeichnen, Werte ablesen
        \item Noch keine Proportionalität/lineare Funktionen
    \end{itemize}
    \item Bekannte Methoden: Einzelarbeit, Partnerarbeit, Plenum, Simulation
\end{itemize}

\subsection{Erfahrungen mit der Lerngruppe}
\begin{itemize}
    \item Bisherige Unterrichtserfahrungen: SuS arbeiten motiviert mit Simulationen, praktische Messungen (100m-Lauf) gut angenommen
    \item Bewährte Sozialformen: Think-Pair-Share, Partnerarbeit beim Zeichnen
    \item Herausforderungen: Achsenbeschriftung, Skalierung, Einheitenkonsistenz
\end{itemize}

% ============================================================================
% 2. SACHANALYSE
% ============================================================================
\section{Sachanalyse}

\subsection{Fachwissenschaftliche Einordnung}
Das \textbf{Weg-Zeit-Diagramm} (s(t)-Diagramm) ist eine zentrale graphische Darstellungsform in der Kinematik. Es visualisiert den funktionalen Zusammenhang zwischen dem zurückgelegten Weg $s$ und der verstrichenen Zeit $t$.

\textbf{Mathematische Grundlage:}
\begin{itemize}
    \item Bei \textbf{gleichförmiger Bewegung} gilt: $s = v \cdot t$ (Ursprungsgerade durch (0|0))
    \item Die \textbf{Steigung} der Geraden entspricht der Geschwindigkeit: $v = \frac{\Delta s}{\Delta t}$
    \item Je steiler die Gerade, desto größer die Geschwindigkeit
    \item Eine \textbf{horizontale Linie} bedeutet Stillstand ($v = 0$)
\end{itemize}

\textbf{Konventionen:}
\begin{itemize}
    \item Zeitachse ($t$): horizontal (x-Achse), Einheit meist Sekunden (s)
    \item Wegachse ($s$): vertikal (y-Achse), Einheit meist Meter (m)
    \item Startpunkt bei $(0|0)$ oder mit Anfangsweg $s_0$
\end{itemize}

\subsection{Kontrastive Betrachtung}
\textbf{Typische Fehlvorstellungen bei Schülern:}
\begin{itemize}
    \item \textbf{Graph = Abbild der Bewegung:} SuS interpretieren das Diagramm als ``Landkarte'' der Bewegung (z.B. steigende Gerade = bergauf fahren)
    \item \textbf{Verwechslung von v(t) und s(t):} Geschwindigkeit vs. Weg
    \item \textbf{Achsenverwechslung:} Zeit auf y-Achse, Weg auf x-Achse
    \item \textbf{Einheitenfehler:} Inkonsistente Einheiten (m und km gemischt)
\end{itemize}

\textbf{Positiver Transfer:}
\begin{itemize}
    \item Koordinatensystem aus Mathematik bekannt
    \item Tabellen aus Stunde 2 (Geschwindigkeit messen) als Datengrundlage
\end{itemize}

\subsection{Kernkonzepte der Stunde}
\begin{enumerate}
    \item \textbf{s(t)-Diagramm als graphische Darstellung:} Zeit auf x-Achse, Weg auf y-Achse
    \item \textbf{Ablesen von Werten:} Zu einem Zeitpunkt $t$ den zugehörigen Weg $s$ ablesen
    \item \textbf{Zeichnen eines Diagramms:} Aus Wertetabelle Punkte einzeichnen und verbinden
    \item \textbf{Qualitative Interpretation:} Steile Gerade = schnell, flache Gerade = langsam, horizontal = Stillstand
    \item \textbf{Vergleich zweier Bewegungen:} Wer ist schneller? Wann überholt einer den anderen?
\end{enumerate}

\subsection{Weiterführende Aspekte}
\begin{itemize}
    \item Anknüpfung Kl. 9/10: Quantitative Steigungsberechnung $v = \frac{\Delta s}{\Delta t}$
    \item Anknüpfung Kl. 9/10: v(t)-Diagramm als Weiterentwicklung
    \item Anknüpfung Kl. 10: Beschleunigte Bewegung (gekrümmte Graphen)
\end{itemize}

% ============================================================================
% 3. DIDAKTISCHE ANALYSE
% ============================================================================
\section{Didaktische Analyse}

\subsection{Stellung der Stunde in der Unterrichtseinheit}
Dies ist die \textbf{3. Stunde} der Unterrichtsreihe ``\reihenthema'' (4 Stunden gesamt).

\begin{tabular}{|c|l|l|}
\hline
\textbf{Std.} & \textbf{Thema} & \textbf{Schwerpunkt} \\
\hline
1 & Bewegungsarten erforschen & Klassifikation (Bahnform, Bewegungsart) \\
\hline
2 & 100m-Lauf -- Geschwindigkeit messen & $v = \frac{s}{t}$, Messung, Berechnung \\
\hline
\textbf{3} & \textbf{s(t)-Diagramme verstehen und zeichnen} & \textbf{Graphische Darstellung} \\
\hline
4 & Bewegungen vergleichen (Stundenleistung) & Anwendung, Test \\
\hline
\end{tabular}

\subsection{Ziele der Unterrichtseinheit}
Am Ende der Reihe können die SuS...
\begin{itemize}
    \item Bewegungen nach Bahnform und Bewegungsart klassifizieren
    \item Die Geschwindigkeit aus Weg und Zeit berechnen
    \item Bewegungen im s(t)-Diagramm darstellen und interpretieren
    \item Bewegungen verschiedener Körper vergleichen
\end{itemize}

\subsection{Legitimation}
\textbf{Rahmenplanbezug (M-V, Orientierungsstufe Kl. 6):}
\begin{itemize}
    \item \textbf{Themenfeld:} Bewegungen -- Geradlinige, gleichförmige Bewegung (~4 USt)
    \item \textbf{Verbindlicher Inhalt:} ``s(t)-Diagramme -- Graphische Darstellung möglich''
    \item \textbf{Kompetenzbereich [E]:} ``Ergebnisse in Texten, Skizzen, Tabellen, \textbf{Diagrammen} dokumentieren''
    \item \textbf{Kompetenzbereich [K]:} ``Arbeitsergebnisse dokumentieren und präsentieren''
\end{itemize}

\textbf{Bildungswert (Klafki):}
\begin{itemize}
    \item \textbf{Gegenwartsbedeutung:} Navigation (Maps), Sportverfolgung, Fahrplan lesen
    \item \textbf{Zukunftsbedeutung:} Graphische Darstellungen in Beruf und Alltag
    \item \textbf{Exemplarische Bedeutung:} Funktionale Zusammenhänge visualisieren
\end{itemize}

\subsection{Didaktische Reduktion}
\textbf{Bewusst ausgeklammert:}
\begin{itemize}
    \item Quantitative Steigungsberechnung (erst Kl. 9)
    \item Negative Wege (Rückwärtsbewegung)
    \item Beschleunigte Bewegung (gekrümmte Graphen, erst Kl. 10)
    \item v(t)-Diagramme
\end{itemize}

\textbf{Begründung:} Fokus auf qualitatives Verständnis und korrektes Zeichnen. Mathematische Voraussetzungen (Steigung, lineare Funktionen) fehlen noch.

\subsection{Strukturierung des Unterrichts}
\begin{enumerate}
    \item \textbf{Einstieg (5 Min):} Problemstellung -- ``Wie können wir Bewegungen sichtbar machen?'' Bezug zu Stunde 2 (Wertetabelle)
    \item \textbf{Erarbeitung (15 Min):} Einführung s(t)-Diagramm anhand eines Beispiels, gemeinsames Zeichnen
    \item \textbf{Übung (15 Min):} Differenzierte Aufgaben -- Ablesen, Zeichnen, Interpretieren
    \item \textbf{Sicherung (10 Min):} Merksätze, Vergleich Tafelbild, Übertrag ins Arbeitsblatt
\end{enumerate}

% ============================================================================
% 4. STUNDENZIELE
% ============================================================================
\section{Stundenziele}

\subsection{Grobziele}
Die SuS erweitern ihre \textbf{Erkenntnisgewinnungskompetenz [E]} und \textbf{Kommunikationskompetenz [K]}, indem sie lernen, Bewegungen graphisch im Weg-Zeit-Diagramm darzustellen und zu interpretieren.

\textbf{Kompetenzbereiche:}
\begin{itemize}
    \item \textbf{Fachkompetenz [S]:} Zusammenhang zwischen Weg und Zeit graphisch erfassen
    \item \textbf{Erkenntnisgewinnung [E]:} Messdaten in Diagrammen dokumentieren
    \item \textbf{Kommunikation [K]:} Fachtypische Darstellungen (Diagramme) verwenden
    \item \textbf{Bewertung [B]:} Bewegungen qualitativ vergleichen
\end{itemize}

\subsection{Feinziele (operationalisiert)}
\begin{tabular}{|c|p{8.5cm}|c|c|}
\hline
\textbf{Nr.} & \textbf{Die SuS können...} & \textbf{AFB} & \textbf{Niveau} \\
\hline
FZ1 & die Achsen eines s(t)-Diagramms korrekt \textbf{beschriften} (t auf x-Achse, s auf y-Achse, Einheiten). & I & alle \\
\hline
FZ2 & aus einem s(t)-Diagramm Werte für Weg und Zeit \textbf{ablesen}. & I & alle \\
\hline
FZ3 & aus einer Wertetabelle ein s(t)-Diagramm \textbf{zeichnen}, indem sie Punkte einzeichnen und verbinden. & II & alle \\
\hline
FZ4 & qualitativ \textbf{erklären}, was eine steile bzw. flache Gerade bedeutet (schnell/langsam). & II & **/*** \\
\hline
FZ5 & zwei Bewegungen im selben Diagramm \textbf{vergleichen} und Aussagen über Geschwindigkeitsunterschiede treffen. & II & **/*** \\
\hline
FZ6 & zu einem gegebenen s(t)-Diagramm eine passende Bewegungsgeschichte \textbf{entwickeln}. & III & *** \\
\hline
\end{tabular}

\vspace{1em}
\textbf{AFB-Verteilung:}
\begin{itemize}
    \item AFB I (Reproduktion): ca. 30\% (FZ1, FZ2)
    \item AFB II (Reorganisation/Transfer): ca. 50\% (FZ3, FZ4, FZ5)
    \item AFB III (Reflexion/Problemlösung): ca. 20\% (FZ6)
\end{itemize}

% ============================================================================
% 5. METHODISCHE ANALYSE
% ============================================================================
\section{Methodische Analyse}

\subsection{Einstieg}
\begin{itemize}
    \item \textbf{Methode:} Problemimpuls mit Rückbezug -- ``Letzte Stunde habt ihr gemessen, wie schnell ihr lauft. Die Ergebnisse stehen in einer Tabelle. Wie können wir das übersichtlicher darstellen?''
    \item \textbf{Begründung:} Knüpft an Vorwissen an, erzeugt kognitive Aktivierung, zeigt Notwendigkeit der Visualisierung
    \item \textbf{Alternative:} Video einer Bewegung mit eingeblendetem Live-Diagramm
\end{itemize}

\subsection{Erarbeitung}
\begin{itemize}
    \item \textbf{Methode:} Lehrergeleitete Einführung am Beispiel eines Fußgängers (Wertetabelle $\to$ Diagramm)
    \item \textbf{Begründung:} Neues Konzept erfordert strukturierte Einführung, Modellierung durch Lehrkraft
    \item \textbf{Scaffolding:}
    \begin{enumerate}
        \item Achsen einzeichnen und beschriften (gemeinsam)
        \item Skalierung festlegen (gemeinsam)
        \item Ersten Punkt einzeichnen (Lehrkraft)
        \item Weitere Punkte einzeichnen (SuS)
        \item Punkte verbinden (SuS)
    \end{enumerate}
    \item \textbf{Alternative:} Entdeckendes Lernen mit Simulation (mehr Zeit nötig)
\end{itemize}

\subsection{Übungsphasen}
\begin{itemize}
    \item \textbf{Methode:} Differenzierte Einzelarbeit mit Arbeitsblatt, anschließend Partnerkontrolle
    \item \textbf{Progression:}
    \begin{itemize}
        \item Aufgabe 1: Werte ablesen (reproduktiv)
        \item Aufgabe 2: Diagramm zeichnen (konstruktiv)
        \item Aufgabe 3: Bewegungen vergleichen (analytisch)
        \item Aufgabe 4: Geschichte erfinden (generativ)
    \end{itemize}
    \item \textbf{Begründung:} Interleaved practice, von gelenkt zu frei
\end{itemize}

\subsection{Sozialformen}
\begin{tabular}{|l|l|l|}
\hline
\textbf{Phase} & \textbf{Sozialform} & \textbf{Begründung} \\
\hline
Einstieg & Plenum & Gemeinsamer Fokus, Aktivierung des Vorwissens \\
\hline
Erarbeitung & Plenum + LSG & Strukturierte Einführung, Modellierung \\
\hline
Übung & Einzelarbeit & Individuelle Verarbeitung, Selbstständigkeit \\
\hline
Kontrolle & Partnerarbeit & Peer-Feedback, Fehlerkorrektur \\
\hline
Sicherung & Plenum & Ergebnissicherung, Klärung offener Fragen \\
\hline
\end{tabular}

\subsection{Medien}
\begin{itemize}
    \item \textbf{Tafel/Whiteboard:} Beispiel-Diagramm, Tafelbild mit Merksätzen
    \item \textbf{Arbeitsblatt:} AB-003-st-diagramme.pdf (mit vorgezeichnetem Koordinatensystem)
    \item \textbf{Digitaler Lernpfad:} LP-003-st-diagramme.html (optional, für Vertiefung/Hausaufgabe)
    \item \textbf{Simulation:} Interaktives s(t)-Diagramm (für Demonstration oder Vertiefung)
\end{itemize}

\subsection{Differenzierung}
\begin{tabular}{|l|p{3.3cm}|p{3.3cm}|p{3.3cm}|}
\hline
& \textbf{Basis (*)} & \textbf{Standard (**)} & \textbf{Erweitert (***)} \\
\hline
\textbf{Inhalt} & Achsen beschriften, Werte ablesen, vorgezeichnete Punkte verbinden & Punkte einzeichnen und verbinden, qualitative Interpretation & Selbst skalieren, Bewegungen vergleichen, Geschichten erfinden \\
\hline
\textbf{Hilfen} & Koordinatensystem vorgegeben, Punkte teilweise eingezeichnet, Wortliste & Koordinatensystem vorgegeben, Tipps zur Skalierung & Nur leeres Raster \\
\hline
\textbf{Produkt} & Beschriftetes Diagramm, abgelesene Werte & Korrekt gezeichnetes Diagramm mit Interpretation & Vergleich, eigene Geschichte \\
\hline
\end{tabular}

\subsection{Erwartungshorizonte}
\begin{itemize}
    \item \textbf{Minimum (alle SuS):} Achsen korrekt beschriften, Werte aus Diagramm ablesen können
    \item \textbf{Regelstandard (Mehrheit):} Aus Wertetabelle ein Diagramm zeichnen, qualitativ interpretieren
    \item \textbf{Optimal (leistungsstarke SuS):} Zwei Bewegungen vergleichen, eigene Bewegungsgeschichte entwickeln
\end{itemize}

\subsection{Überprüfung der Teilziele}
\begin{tabular}{|c|l|l|}
\hline
\textbf{FZ} & \textbf{Überprüfung durch} & \textbf{Zeitpunkt} \\
\hline
FZ1 & Tafelbild-Vergleich, AB Kasten 1 & Erarbeitung, Sicherung \\
\hline
FZ2 & AB Aufgabe 1 (Ablesen) & Übungsphase \\
\hline
FZ3 & AB Aufgabe 2 (Zeichnen) & Übungsphase \\
\hline
FZ4 & Mündliche Beiträge, AB Aufgabe 3 & Sicherung, Übungsphase \\
\hline
FZ5 & AB Aufgabe 3 (Vergleich) & Übungsphase \\
\hline
FZ6 & AB Aufgabe 4 (***) & Übungsphase/Hausaufgabe \\
\hline
\end{tabular}

% ============================================================================
% 6. VERLAUFSPLANUNG
% ============================================================================
\section{Verlaufsplanung}

\begin{longtable}{|p{0.8cm}|p{1.4cm}|p{4.2cm}|p{3cm}|p{1.5cm}|p{1.8cm}|}
\hline
\textbf{Zeit} & \textbf{Phase} & \textbf{Inhalt / Lehrerhandeln} & \textbf{Schülerhandeln} & \textbf{SF} & \textbf{Medien} \\
\hline
\endhead

0--5 & Einstieg & L zeigt Wertetabelle aus Stunde 2 (100m-Lauf). ``Wie können wir diese Daten übersichtlicher darstellen?'' Sammelt Vorschläge. & SuS erinnern sich an Tabelle, nennen Ideen (Diagramm, Grafik) & Plenum & Beamer/ Tafel \\
\hline

5--10 & Erarbei-tung I & L erklärt Grundaufbau: ``Das nennt man Weg-Zeit-Diagramm.'' Zeichnet Koordinatensystem an Tafel, beschriftet Achsen (t, s) mit Einheiten. & SuS beobachten, übernehmen Beschriftungsregel in AB Kasten 1 & Plenum/ LSG & Tafel, AB \\
\hline

10--18 & Erarbei-tung II & L zeigt Beispiel-Wertetabelle (Fußgänger: 0s/0m, 2s/10m, 4s/20m, 6s/30m). Demonstriert Skalierung, trägt Punkte ein, verbindet sie. Bespricht: ``Was bedeutet die Gerade?'' & SuS verfolgen Demonstration, beantworten Fragen, notieren Merksatz (Kasten 2) & Plenum/ LSG & Tafel, AB \\
\hline

18--30 & Übung & L verteilt AB (falls nicht schon geschehen). ``Bearbeitet jetzt die Aufgaben 1--3 in eurem Tempo.'' Geht herum, gibt individuelle Hilfe. & SuS bearbeiten differenzierte Aufgaben: Ablesen (1), Zeichnen (2), Vergleichen (3) & EA & AB, Lineal \\
\hline

30--35 & Partner-kontrolle & ``Vergleicht eure Ergebnisse mit eurem Nachbarn. Besprecht Unterschiede.'' & SuS vergleichen Diagramme, diskutieren Fehler, korrigieren & PA & AB \\
\hline

35--42 & Siche-rung & L bespricht Lösungen im Plenum (Dokumentenkamera/Tafel). Klärt typische Fehler (Achsenbeschriftung, Skalierung). Wiederholt Merksätze. & SuS korrigieren, stellen Fragen, vervollständigen AB & Plenum & Doku-Kamera/ Tafel \\
\hline

42--45 & Abschluss & L gibt Ausblick: ``Nächste Stunde wendet ihr das in einem Test an.'' Benennt HA (optional: Aufgabe 4 oder Lernpfad). & SuS notieren HA, räumen auf & Plenum & -- \\
\hline
\end{longtable}

\textbf{Legende:} EA = Einzelarbeit, PA = Partnerarbeit, LSG = Lehrer-Schüler-Gespräch

% ============================================================================
% 7. REFLEXION (nach der Durchführung)
% ============================================================================
\section{Reflexion (nach der Durchführung)}

\textit{Dieser Abschnitt wird nach der Unterrichtsdurchführung ausgefüllt.}

\subsection{Realisierung der Lernziele}
\begin{tabular}{|c|c|l|}
\hline
\textbf{Feinziel} & \textbf{Erreicht?} & \textbf{Evidenz/Anmerkung} \\
\hline
FZ1 & $\Box$ ja / $\Box$ teilweise / $\Box$ nein & \\
\hline
FZ2 & $\Box$ ja / $\Box$ teilweise / $\Box$ nein & \\
\hline
FZ3 & $\Box$ ja / $\Box$ teilweise / $\Box$ nein & \\
\hline
FZ4 & $\Box$ ja / $\Box$ teilweise / $\Box$ nein & \\
\hline
FZ5 & $\Box$ ja / $\Box$ teilweise / $\Box$ nein & \\
\hline
FZ6 & $\Box$ ja / $\Box$ teilweise / $\Box$ nein & \\
\hline
\end{tabular}

\subsection{Wirksamkeit des Vorgehens}
\begin{itemize}
    \item Was hat gut funktioniert? ...
    \item Was hat nicht funktioniert? ...
    \item Unvorhergesehene Ereignisse: ...
\end{itemize}

\subsection{Konsequenzen für die Weiterführung}
\begin{itemize}
    \item Für die nächste Stunde (Stundenleistung): ...
    \item Für ähnliche Stunden: ...
\end{itemize}

% ============================================================================
% 8. ANLAGEN
% ============================================================================
\section{Anlagen}

\begin{enumerate}
    \item Arbeitsblatt (AB-003-st-diagramme.pdf)
    \item Musterlösung (ML-003-st-diagramme.pdf)
    \item Lernpfad (LP-003-st-diagramme.html)
    \item Lehrerhinweise (LH-003-st-diagramme.pdf)
    \item Tafelbild (Abschnitt \ref{sec:tafelbild})
\end{enumerate}

\subsection{Tafelbild}\label{sec:tafelbild}

\begin{center}
\fbox{\parbox{0.9\textwidth}{
\textbf{\large Weg-Zeit-Diagramm (s(t)-Diagramm)}

\vspace{1em}

\begin{minipage}{0.45\textwidth}
\textbf{Wertetabelle:}

\begin{tabular}{|c|c|}
\hline
$t$ in s & $s$ in m \\
\hline
0 & 0 \\
2 & 10 \\
4 & 20 \\
6 & 30 \\
\hline
\end{tabular}
\end{minipage}
\hfill
\begin{minipage}{0.45\textwidth}
\textbf{Diagramm:}

\begin{tikzpicture}[scale=0.6]
    \draw[->] (0,0) -- (7,0) node[right] {$t$ in s};
    \draw[->] (0,0) -- (0,4.5) node[above] {$s$ in m};
    \foreach \x in {0,2,4,6} \draw (\x,0.1) -- (\x,-0.1) node[below] {\x};
    \foreach \y/\label in {0/0,1/10,2/20,3/30} \draw (0.1,\y) -- (-0.1,\y) node[left] {\label};
    \draw[thick,blue] (0,0) -- (2,1) -- (4,2) -- (6,3);
    \fill[blue] (0,0) circle (3pt);
    \fill[blue] (2,1) circle (3pt);
    \fill[blue] (4,2) circle (3pt);
    \fill[blue] (6,3) circle (3pt);
\end{tikzpicture}
\end{minipage}

\vspace{1em}

\textbf{Merksätze:}
\begin{itemize}
    \item Zeitachse ($t$) immer \textbf{waagerecht} (x-Achse)
    \item Wegachse ($s$) immer \textbf{senkrecht} (y-Achse)
    \item Steile Gerade $\to$ \textbf{schnelle} Bewegung
    \item Flache Gerade $\to$ \textbf{langsame} Bewegung
    \item Waagerechte Linie $\to$ \textbf{Stillstand}
\end{itemize}
}}
\end{center}

% ============================================================================
% 9. LÖSUNGEN
% ============================================================================
\section{Lösungen}

\subsection{Lösungen Arbeitsblatt}

\textbf{Kasten 1 -- Achsenbeschriftung:}
\begin{itemize}
    \item x-Achse (waagerecht): Zeit $t$ in Sekunden (s)
    \item y-Achse (senkrecht): Weg $s$ in Metern (m)
\end{itemize}

\textbf{Kasten 2 -- Merksatz:}
Im s(t)-Diagramm zeigt eine \textbf{steile Gerade} eine \textbf{schnelle} Bewegung. Eine \textbf{flache Gerade} zeigt eine \textbf{langsame} Bewegung. Eine \textbf{waagerechte Linie} bedeutet \textbf{Stillstand}.

\textbf{Aufgabe 1 -- Werte ablesen (*):}
\begin{enumerate}[label=\alph*)]
    \item Nach $t = 4$ s hat der Radfahrer $s = 40$ m zurückgelegt.
    \item Nach $t = 10$ s hat der Radfahrer $s = 100$ m zurückgelegt.
    \item Für $s = 60$ m benötigt der Radfahrer $t = 6$ s.
\end{enumerate}

\textbf{Aufgabe 2 -- Diagramm zeichnen (**):}

Wertetabelle Jogger: (0s|0m), (5s|25m), (10s|50m), (15s|75m), (20s|100m)

$\to$ Gerade durch Ursprung mit Steigung 5 m/s

\textbf{Aufgabe 3 -- Bewegungen vergleichen (**):}
\begin{enumerate}[label=\alph*)]
    \item Der Radfahrer ist schneller, weil seine Gerade steiler ist (mehr Weg in gleicher Zeit).
    \item Nach $t = 10$ s: Radfahrer bei 100 m, Jogger bei 50 m $\to$ Radfahrer ist 50 m voraus.
    \item Sie treffen sich nicht, da der Radfahrer immer schneller ist und beide bei $s = 0$ starten.
\end{enumerate}

\textbf{Aufgabe 4 -- Bewegungsgeschichte (***):}

Beispiel für ein Diagramm mit Stillstand:

``Lisa fährt mit dem Fahrrad zur Schule. In den ersten 2 Minuten fährt sie 400 m. Dann muss sie an einer roten Ampel warten (1 Minute Stillstand). Danach fährt sie noch 3 Minuten und legt dabei 600 m zurück.''

\subsection{Erwartete Schülerantworten}

\textbf{Frage: ``Was bedeutet eine steile Gerade?''}
\begin{itemize}
    \item Richtig: ``Der Körper bewegt sich schnell.'' / ``Er legt viel Weg in kurzer Zeit zurück.''
    \item Typischer Fehler: ``Der Körper fährt bergauf.'' (Graph = Abbild-Fehlvorstellung)
\end{itemize}

\textbf{Frage: ``Was bedeutet eine waagerechte Linie?''}
\begin{itemize}
    \item Richtig: ``Der Körper steht still.'' / ``Der Weg ändert sich nicht.''
\end{itemize}

% ============================================================================
% 10. LITERATUR
% ============================================================================
\section{Literatur}

\subsection{Rahmenplan}
\begin{itemize}
    \item Ministerium für Bildung und Kindertagesförderung M-V / IQ M-V (2022): Rahmenplan Physik -- Orientierungsstufe Klasse 6. Erprobungsfassung.
\end{itemize}

\subsection{Fachdidaktische Literatur}
\begin{itemize}
    \item Kircher, E., Girwidz, R., Häußler, P. (2020): Physikdidaktik. Springer.
    \item Wiesner, H. et al. (2011): Schülervorstellungen in der Physik. Aulis.
\end{itemize}

\subsection{Digitale Quellen}
\begin{itemize}
    \item LEIFI-Physik: Gleichförmige Bewegung. \url{https://www.leifiphysik.de/}
    \item PhET Simulations: Moving Man. \url{https://phet.colorado.edu/}
\end{itemize}

% ============================================================================
% 11. DIDAKTIK-SCORE
% ============================================================================
\newpage
\section{Didaktik-Score}

\begin{scorebox}
\textbf{Bewertung der didaktischen Qualität nach 10 Dimensionen}\\
\textbf{Ziel-Score: $\geq$ 9.0 (Premium-Qualität)}

\vspace{1em}
\begin{tabular}{|c|l|c|c|p{4.5cm}|}
\hline
\textbf{\#} & \textbf{Dimension} & \textbf{Gew.} & \textbf{Score} & \textbf{Notiz} \\
\hline
1 & Differenzierung & 15\% & 9/10 & 3 qualitativ verschiedene Niveaus; Hilfen gestuft \\
\hline
2 & Taxonomie (AFB) & 10\% & 9/10 & 30/50/20 -- leicht mehr AFB I, angemessen für Kl. 6 \\
\hline
3 & Lernziel-Alignment & 10\% & 10/10 & Direkt aus RP-MV: ``s(t)-Diagramme'' \\
\hline
4 & Aktivierung & 10\% & 9/10 & Zeichnen = konstruktiv; Geschichte = generativ \\
\hline
5 & Scaffolding & 10\% & 9/10 & Gestufte Demo, vorstrukturierte AB-Varianten \\
\hline
6 & Feedback-Qualität & 10\% & 8/10 & Partnerkontrolle + Plenumsbesprechung; LP mit Auto-Feedback ergänzen \\
\hline
7 & Sprachliche Klarheit & 10\% & 9/10 & Altersgerecht (Kl. 6), Fachbegriffe eingeführt \\
\hline
8 & Motivation \& Relevanz & 10\% & 9/10 & Bezug zu 100m-Lauf (eigene Daten), Alltagsbeispiele \\
\hline
9 & Inklusion (UDL) & 10\% & 8/10 & Visuell (Diagramm) + verbal; LRS-Hilfen genannt \\
\hline
10 & Praktikabilität & 5\% & 10/10 & 45 Min realistisch, Standardmaterial \\
\hline
\multicolumn{3}{|r|}{\textbf{GESAMT:}} & \textbf{9.0/10} & \textcolor{successgreen}{\textbf{$\checkmark$ FREIGABE}} \\
\hline
\end{tabular}

\vspace{1em}
\textbf{Bewertungsschwellen:}
\begin{itemize}
    \item[\checkmark] \textcolor{successgreen}{$\geq$ 9.0: \textbf{FREIGABE} (Premium-Qualität)}
    \item[$\Box$] \textcolor{warningorange}{8.0--8.9: Iteration empfohlen}
    \item[$\Box$] 6.0--7.9: Überarbeitung erforderlich
    \item[$\Box$] $<$ 6.0: Grundlegende Neukonzeption
\end{itemize}
\end{scorebox}

\subsection{Score-Begründungen}

\textbf{1. Differenzierung (15\%) -- 9/10}
\begin{itemize}
    \item Drei qualitativ verschiedene Niveaus: (*) Ablesen + vorgefertigte Punkte, (**) selbst zeichnen + interpretieren, (***) vergleichen + generieren
    \item Nicht nur mehr/weniger Aufgaben, sondern andere kognitive Anforderungen
    \item Hilfen gestuft (vorstrukturiertes Koordinatensystem vs. leeres Raster)
\end{itemize}

\textbf{2. Taxonomie/AFB (10\%) -- 9/10}
\begin{itemize}
    \item Verteilung 30/50/20 -- leicht höherer AFB-I-Anteil für Kl. 6 angemessen
    \item Operatoren korrekt verwendet (beschriften, ablesen, zeichnen, erklären, vergleichen, entwickeln)
\end{itemize}

\textbf{3. Lernziel-Alignment (10\%) -- 10/10}
\begin{itemize}
    \item Direkt aus Rahmenplan MV Klasse 6: ``s(t)-Diagramme -- Graphische Darstellung möglich''
    \item Kompetenzbereiche [E] und [K] explizit adressiert
\end{itemize}

\textbf{4. Aktivierung (10\%) -- 9/10}
\begin{itemize}
    \item Zeichnen = konstruktive Aktivität (nicht nur ankreuzen)
    \item Bewegungsgeschichte erfinden = generative Aktivität
    \item Einstieg aktiviert Vorwissen
\end{itemize}

\textbf{5. Scaffolding (10\%) -- 9/10}
\begin{itemize}
    \item Gestufte Demonstration (gemeinsam $\to$ selbstständig)
    \item Vorstrukturierte AB-Varianten für schwächere SuS
    \item Lösung wird nicht verraten, sondern schrittweise erarbeitet
\end{itemize}

\textbf{6. Feedback-Qualität (10\%) -- 8/10}
\begin{itemize}
    \item Partnerkontrolle ermöglicht Peer-Feedback
    \item Plenumsbesprechung klärt typische Fehler
    \item \textit{Verbesserungspotential:} Digitaler Lernpfad mit automatischem Feedback ergänzen
\end{itemize}

\textbf{7. Sprachliche Klarheit (10\%) -- 9/10}
\begin{itemize}
    \item Altersgerechte Sprache für Klasse 6
    \item Neue Fachbegriffe (Weg-Zeit-Diagramm, Achsen) werden eingeführt und erklärt
    \item Kurze, klare Arbeitsanweisungen
\end{itemize}

\textbf{8. Motivation \& Relevanz (10\%) -- 9/10}
\begin{itemize}
    \item Bezug zu eigenen Messdaten (100m-Lauf aus Stunde 2)
    \item Alltagsbeispiele (Fußgänger, Radfahrer, Jogger)
    \item ARCS: Attention (eigene Daten), Relevance (Sport, Verkehr), Confidence (gestufte Hilfen)
\end{itemize}

\textbf{9. Inklusion/UDL (10\%) -- 8/10}
\begin{itemize}
    \item Visuell (Diagramm) + verbal (Merksätze) + handlungsorientiert (Zeichnen)
    \item LRS-Hilfen explizit genannt (vergrößerte Schrift, mehr Zeit)
    \item \textit{Verbesserungspotential:} Zusätzliche Wortliste für DaZ-SuS
\end{itemize}

\textbf{10. Praktikabilität (5\%) -- 10/10}
\begin{itemize}
    \item 45 Minuten realistisch planbar
    \item Standardmaterial (Tafel, AB, Lineal)
    \item Keine aufwendige Vorbereitung nötig
\end{itemize}

\subsection{Verbesserungsmaßnahmen}
Der Score von 9.0 erreicht die Freigabe-Schwelle. Optionale Verbesserungen für zukünftige Iterationen:

\begin{enumerate}
    \item \textbf{Feedback-Qualität (+0.5):} Digitalen Lernpfad mit automatischer Auswertung der Diagramme entwickeln
    \item \textbf{Inklusion (+0.3):} Wortliste mit Fachbegriffen (deutsch/türkisch/arabisch) für DaZ-SuS erstellen
    \item \textbf{Aktivierung (+0.2):} Optional: Eigene Bewegungsdaten (Schulweg) als Grundlage für Diagramm nutzen
\end{enumerate}

% ============================================================================
% ENDE
% ============================================================================
\end{document}
